problem:  
Let \( \mathcal{L}(n) \) denote the real algebraic variety of \(n\)-dimensional Lie brackets on \(\mathbb{R}^n\), and let \(O(n)\) act on \(\mathcal{L}(n)\) by change of orthonormal basis.  Lauret’s results identify **Ricci soliton left‑invariant metrics** with zeros of a moment‑map‑type functional (“the bracket flow moment map”).  Each **homogeneous Ricci soliton** corresponds to an \(O(n)\)-orbit representing an **algebraic soliton** point
\[
[\mu]\in \mathcal{L}(n)/O(n),\qquad \text{satisfying } \operatorname{Ric}_\mu = c\,I + S(D_\mu)
\]
for some derivation \(D_\mu\) of the Lie algebra \((\mathbb{R}^n,\mu)\).  

**Question:**  
For the variety \( \mathcal{L}(n) \) equipped with the moment map structure arising from the Ricci operator, consider the algebraic soliton points \([\mu]\) satisfying the soliton equation.  

1.  Do the **closures of distinct \(O(n)\)-orbits** of algebraic soliton brackets ever meet?  Equivalently, can two non‑isomorphic solvmanifolds with distinct bracket types lie in the same **GIT degeneration class** under the bracket flow while both being algebraic solitons?  
2.  If such intersections occur, do they correspond to geometric limits of homogeneous Ricci solitons—i.e., can the bracket flow degenerate one algebraic soliton into a *non‑isomorphic* algebraic limit preserving the soliton condition?  
3.  Is the set of algebraic soliton orbits \( \mathcal{S}_n \subset \mathcal{L}(n)/O(n) \) **closed** under degeneration (so every geometric limit of algebraic solitons is algebraic), or does it exhibit a frontier of non‑algebraic brackets?  

Why is it a "good" problem:  
This formulation shifts the focus from uniqueness to **moduli and boundary behavior** of algebraic solitons within the global geometry of the Lie‑bracket variety.  It builds directly on Lauret’s moment‑map picture and bracket‑flow framework, but the questions above remain entirely unresolved: the structure of degenerations between algebraic soliton orbits, and whether the class of algebraic solitons is closed under limits, are genuinely open.  Resolving these would connect geometric invariant theory, the dynamics of the homogeneous Ricci flow, and the stratification of the space of left‑invariant metrics.  The statement is algebraically simple (do algebraic solitons form a closed and isolated orbit family under \(O(n)\)?), yet its resolution would profoundly clarify the global geometry of homogeneous Ricci solitons, achieving the “narrow entry, deep consequence” ideal of a top‑level problem.